\chapter{Docker Workflow}

\section{Docker Image}

The official OmniLaTeX Docker image is published to GitHub Container
Registry:

\begin{minted}{text}
ghcr.io/wyattau/omnilatex-docker
\end{minted}

The image bundles TeX Live 2026 with all template dependencies pre-installed.
Font caches are pre-warmed, eliminating the cold-start penalty on first
compilation.  The compressed image is approximately 2\,GB.

\section{Pulling the Image}

\begin{minted}{bash}
docker pull ghcr.io/wyattau/omnilatex-docker:latest
\end{minted}

Pin to a specific version for reproducibility:

\begin{minted}{bash}
docker pull ghcr.io/wyattau/omnilatex-docker:v2026.1.0
\end{minted}

\section{Build Commands}

A full production build inside Docker:

\begin{minted}{bash}
docker run --rm \
  --entrypoint "" \
  -v "$PWD:/workspace" \
  -w /workspace \
  ghcr.io/wyattau/omnilatex-docker:latest \
  bash -c "BUILD_MODE=prod python build.py"
\end{minted}

Development build:

\begin{minted}{bash}
docker run --rm \
  --entrypoint "" \
  -v "$PWD:/workspace" \
  -w /workspace \
  ghcr.io/wyattau/omnilatex-docker:latest \
  bash -c "BUILD_MODE=dev python build.py"
\end{minted}

Ultra-fast iteration:

\begin{minted}{bash}
docker run --rm \
  --entrypoint "" \
  -v "$PWD:/workspace" \
  -w /workspace \
  ghcr.io/wyattau/omnilatex-docker:latest \
  bash -c "BUILD_MODE=ultra python build.py"
\end{minted}

\section{ENTRYPOINT Bypass}

The Docker image sets \texttt{ENTRYPOINT} to \texttt{latexmk "\$@"}, which
means any arguments passed to \texttt{docker run} are forwarded directly to
\texttt{latexmk}.  To run arbitrary shell commands (such as \texttt{build.py}
or \texttt{make}), you \emph{must} use:

\begin{minted}{text}
--entrypoint ""
\end{minted}

Without this flag, the container attempts to pass your command as arguments
to \texttt{latexmk}, which will fail.

\section{Volume Mounting}

The standard volume-mount pattern maps the current directory into the
container's workspace:

\begin{minted}{bash}
-v "$PWD:/workspace" -w /workspace
\end{minted}

This ensures that output PDFs are written directly to your local filesystem.
No copy step is needed.

\section{Font Cache}

The Docker image ships with pre-warmed font caches for all bundled fonts.
On a cold host, this saves 30--60 seconds on the first \texttt{latexmk}
pass compared to a bare TeX Live installation.

\section{Cross-Platform}

The Docker image runs on all major platforms:

\begin{itemize}
  \item \textbf{Linux} --- native Docker Engine.
  \item \textbf{macOS} --- Docker Desktop (Intel or Apple Silicon).
  \item \textbf{Windows} --- WSL2 with Docker Desktop.
\end{itemize}

No platform-specific configuration is required.  The same
\texttt{docker run} command works identically everywhere.
